\section{Foundations}

% BUDGET: ~60s
\begin{frame}{Inconsistency measure $\mathcal{I}(K)$}
    \begin{block}{Inconsistency measure $\mathcal{I} : \mathbb{K} \to \mathbb{R}_{\geq 0}$ [Thimm, 2019]}
        \begin{itemize}
            \item Higher value means more inconsistency, $\mathcal{I}(K) = 0$ iff $K$ consistent
            \item Rationality postulates such as Monotonicity and Free Formula Independence pin down sensible measures
            \item Contension $\mathcal{I}_c$ [Grant and Hunter, 2011] counts the atoms it must assign the truth value $B$, both true and false
        \end{itemize}
    \end{block}

    \begin{exampleblock}{Contension $\mathcal{I}_c$}
        $K = \{a \land b,\ \neg a,\ \neg b\}$ forces $a$ and $b$ to $B$, so $\mathcal{I}_c(K) = 2$.
    \end{exampleblock}
    \sectioncite{Thesis \S2.1}
\end{frame}

% BUDGET: ~60s
\begin{frame}{Classical planning}
    \small A PDDL Domain and Problem translate to a STRIPS tuple $\Pi = \langle P, O, I, G \rangle$.

    \begin{block}{STRIPS notation $\Pi = \langle P, O, I, G \rangle$ [Bylander 1994]}
        \footnotesize
        \begin{itemize}
            \item $P$ propositions, $O$ operators, $I \subseteq P$ initial state, $G \subseteq P$ goal
            \item Each operator $o = \langle \text{precond}(o), \text{add}(o), \text{delete}(o) \rangle$
            \item Plan is an applicable operator sequence from $I$ reaching a state $\supseteq G$. Unsolvable means none exists
        \end{itemize}
    \end{block}
    \begin{columns}[T]
        \begin{column}{0.6\textwidth}
            \begin{exampleblock}{Rover Mission, the drill subproblem}
                \footnotesize
                $P = \{\text{arm\_deployed},\ \text{drill\_bit},\ \text{hole\_drilled},\ \text{soil\_sampled}\}$ \\[2pt]
                $O = \{\ \text{drill} = \langle \{\text{arm\_deployed}, \text{drill\_bit}\},\ \{\text{hole\_drilled}\},\ \emptyset \rangle,$ \\[1pt]
                $\hphantom{O = \{\ } \text{sample} = \langle \{\text{hole\_drilled}\},\ \{\text{soil\_sampled}\},\ \emptyset \rangle\ \}$ \\[2pt]
                $I = \{\text{arm\_deployed}\}, \qquad G = \{\text{soil\_sampled}\}$
            \end{exampleblock}
        \end{column}
        \begin{column}{0.4\textwidth}
            \vspace*{-0.5em}
            \centering
            \begin{tikzpicture}[
        scale=0.78, transform shape,
        prop/.style={circle, draw, thick, font=\small, align=center, inner sep=1.5pt},
        missing/.style={prop, dashed},
        initial/.style={prop, double},
        goal/.style={prop, fill=black!16},
        oplbl/.style={font=\footnotesize\itshape},
        >={Stealth[length=1.5mm]},
    ]
    \node[initial] (arm)      at (0, 0)      {arm\_\\deployed};
    \node[missing] (drillbit) at (2.8, 0)    {drill\_\\bit};
    \node[prop]    (hole)     at (0, -2.4)   {hole\_\\drilled};
    \node[goal]    (soil)     at (2.8, -2.4) {soil\_\\sampled};
    \draw[->, thick] (arm)      -- (hole);
    \draw[->, thick] (drillbit) -- (hole);
    \node[oplbl] at (0.5, -1.3) {drill};
    \draw[->, thick] (hole) -- node[oplbl, below] {sample} (soil);
\end{tikzpicture}

        \end{column}
    \end{columns}

    \sectioncite{Thesis \S2.2}
\end{frame}

% BUDGET: ~60s
\begin{frame}{Answer Set Programming (ASP)}
    \begin{block}{Execution traces and reasoning modes}
        \begin{itemize}
            \item Each answer set of the encoding is one execution trace, a state sequence from $I$
            \item Brave reasoning unions the answer sets, an atom holds if some trace produces it
            \item Cautious reasoning intersects them, an atom holds only if every trace produces it [Eiter and Gottlob, 1995]
        \end{itemize}
    \end{block}

    \begin{exampleblock}{Brave and cautious over two traces}
        Trace 1 makes $\{a, b\}$ reachable, trace 2 makes $\{a, c\}$ reachable. \\[2pt]
        $\text{BRAVE} = \{a, b, c\}$ reachable by some trace, \quad $\text{CAUTIOUS} = \{a\}$ by every trace.
    \end{exampleblock}

    \sectioncite{Thesis \S2.4}
\end{frame}
